\chapter{Вещественно-структурное кратностное факторпространство блока инцидентности}
\label{ch:v8-real-incidence-multiplicity-quotient-gate}

\section{Вещественная структура меняет ориентацию, а не копийную метку}

Ориентированное поле переноса имеет тип
\begin{equation}
 A:\mathbb C^{11}\longrightarrow\mathbb C^{10}.
 \label{eq:v8-real-oriented-transfer}
\end{equation}
Его партнёр вещественной структуры имеет противоположный тип
\begin{equation}
 A^*:\mathbb C^{10}\longrightarrow\mathbb C^{11}.
 \label{eq:v8-real-opposite-transfer}
\end{equation}
Поэтому оператор вещественной структуры не является антилинейным эндоморфизмом одной
ориентированной половины. Он требует удвоенного носителя
\begin{equation}
 \mathcal T_{15}\oplus\overline{\mathcal T}_{15},
 \qquad
 J(A,\overline B)=(B,\overline A).
 \label{eq:v8-real-transfer-doubling}
\end{equation}
Его неподвижная самосопряжённая часть по-прежнему имеет 30 вещественных
параметров, то есть столько же, сколько исходное комплексное поле. Условие вещественной структуры
удвоение не создаёт дополнительного факторпространства размерности.

\section{Вырожденный блок и его коммутант}

Проблемный совместный блок $(\mathcal C_G,\mathcal B)=(1,0)$ имеет
комплексную размерность восемь и содержит проектор $P_I$ ранга четыре.
Ограничение полного калибровочного представления на этот блок имеет коммутант
\begin{equation}
 \dim_{\mathbb C}\operatorname{End}_G(\mathcal T_{(1,0)})=10.
 \label{eq:v8-real-degenerate-block-commutant}
\end{equation}
Следовательно, неоднозначность шире простого скалярного $U(2)$: существуют
ненулевые операторы, коммутирующие с калибровочным действием и смешивающие инцидентностные и тяжёлые
подпространства.

Из такого эрмитова генератора $K$ построено семейство
\begin{equation}
 P_I(\theta)=e^{i\theta K}P_Ie^{-i\theta K}.
 \label{eq:v8-real-compatible-projector-orbit}
\end{equation}
Для всех проверенных $\theta$ проектор имеет комплексный ранг четыре и
коммутирует с калибровочным действием с остатком меньше $1.5\cdot10^{-15}$.
При $\theta=1.5$ расстояние от исходного $P_I$ равно
$1.36328\ldots$, поэтому орбита нетривиальна.

\section{совместимость с вещественной структурой всей орбиты}

На удвоенном блоке каждый проектор поднимается как
\begin{equation}
 \widehat P(\theta)
 =\operatorname{diag}\bigl(P_I(\theta),\overline{P_I(\theta)}\bigr).
 \label{eq:v8-real-doubled-projector}
\end{equation}
Для обменного оператора вещественной структуры выполнено
\begin{equation}
 J\widehat P(\theta)J^{-1}=\widehat P(\theta)
 \label{eq:v8-real-projector-compatibility}
\end{equation}
с нулевым машинным остатком на всей выборке. На неподвижной относительно вещественной структуры части все
эти проекторы имеют один ранг восемь. Значит, условие вещественной структуры допускает всю
непрерывную орбиту и не выбирает копию инцидентности.

Физический полуслед также сокращает удвоение равномерно:
\begin{equation}
 \frac12\Tr\operatorname{diag}(X,\overline X)^*
                 \operatorname{diag}(X,\overline X)
 =\Tr X^*X.
 \label{eq:v8-real-half-trace-uniform-again}
\end{equation}
Поэтому он не создаёт отношения $m_I/m_H$.

\begin{theorem}[удвоение вещественной структуры не снимает кратность $4+4$]
Ориентационно-обменная вещественная структура переводит модуль переноса в
сопряжённую противоположную половину и сохраняет физическую вещественную
размерность 30. В вырожденном блоке существует непрерывное семейство
калибровочно- и вещественно-структурно совместимых проекторов ранга четыре. Следовательно, условие вещественной структуры
условие не выделяет $P_I$.
\end{theorem}

\begin{nogocriterion}[Запрет фиктивного факторпространства вещественной структуры]
Нельзя отождествлять самосопряжённость относительно вещественной структуры с условием вещественности каждой
матрицы $A$ и тем самым искусственно делить модуль переноса пополам. Нельзя
также выбирать один проектор из орбиты $P_I(\theta)$ только по его совпадению
с желаемым вакуумом: все проекторы проходят одинаковые калибровочные и вещественно-структурные проверки.
\end{nogocriterion}

\begin{protocol}
\begin{enumerate}
 \item ориентированная половина: \textbf{$15$ комплексных измерений};
 \item удвоение вещественной структуры: \textbf{$30$ комплексных измерений};
 \item неподвижная относительно вещественной структуры часть: \textbf{$30$ вещественных измерений};
 \item проблемный блок: \textbf{$8$ комплексных измерений};
 \item инцидентностная и тяжёлая части ранги: \textbf{$4+4$};
 \item калибровочный коммутант блока: \textbf{$10$ комплексных измерений};
 \item непрерывное смешивание копий: \textbf{есть};
 \item совместимость с вещественной структурой всей орбиты: \textbf{да};
 \item вещественный полуслед фиксирует отношение масс: \textbf{нет};
 \item факторпространства вещественной структуры выбирает $P_I$: \textbf{нет};
 \item полный родительское действие: \textbf{не получен};
 \item следующий гейт: \textbf{бимодульный селектор кратности}.
\end{enumerate}
\end{protocol}

Машинный сертификат сохранён в
\path{s2t_v8_real_incidence_multiplicity_quotient_gate_results.json}.